% TOP_PROBLEM: {"id": 3349, "title": "Algebraic independence of pi and e", "category_id": 1, "category": {"id": 1, "name": "number_theory", "display_name": "Number Theory", "description": "Properties of integers, prime numbers, Diophantine equations.", "slug": "number-theory", "order_index": 1, "created_at": "2026-07-31T15:26:25.670Z"}, "status": "open"}
% FIRST_POSTED: 2026-09-27
% !TeX program = lualatex
% Compile twice: lualatex open_problems_500.tex
% All references and prose are embedded; no BibTeX database or external inputs.
\documentclass[11pt,a4paper]{article}
\usepackage[margin=24mm,headheight=14pt]{geometry}
\usepackage{fontspec}
\setmainfont{Latin Modern Roman}
\setsansfont{Latin Modern Sans}
\setmonofont{Latin Modern Mono}
\usepackage{amsmath,amssymb,mathtools}
\usepackage{unicode-math}
\setmathfont{Latin Modern Math}
\usepackage{microtype}
\usepackage{enumitem,xurl,needspace}
\usepackage[dvipsnames]{xcolor}
\usepackage{hyperref}
\hypersetup{unicode=true,colorlinks=true,linkcolor=MidnightBlue,urlcolor=MidnightBlue,
 pdftitle={500 Mathematical Problem Statements},pdfauthor={Source-annotated compilation},bookmarksnumbered=true}
\usepackage{bookmark}
\usepackage{tocloft}
\setlength{\cftsecnumwidth}{3em}
\cftsetpnumwidth{2.5em}
\cftsetrmarg{4em}
\usepackage{titlesec}
\titleformat{\section}{\Large\bfseries\raggedright}{\thesection}{0.7em}{}
\usepackage{fancyhdr}
\pagestyle{fancy}\fancyhf{}
\fancyhead[L]{\small\sffamily Mathematical problem statements}
\fancyhead[R]{\small Source edition 22}
\fancyfoot[C]{\thepage}
\setlength{\parindent}{0pt}
\setlength{\parskip}{5pt plus 1pt}
\setlength{\emergencystretch}{3em}
\setcounter{tocdepth}{1}
\setcounter{secnumdepth}{1}
\setlist[itemize]{leftmargin=3em,itemsep=5pt,topsep=4pt,parsep=0pt}
\allowdisplaybreaks
\newcommand{\N}{\mathbb{N}}
\newcommand{\Z}{\mathbb{Z}}
\newcommand{\Q}{\mathbb{Q}}
\newcommand{\R}{\mathbb{R}}
\newcommand{\C}{\mathbb{C}}
\newcommand{\F}{\mathbb{F}}
\newcommand{\A}{\mathbb{A}}
\newcommand{\PP}{\mathbb P}
\newcommand{\E}{\mathbb{E}}
\newcommand{\eps}{\varepsilon}
\newcommand{\dd}{\,\mathrm d}
\newcommand{\sref}[2]{\hyperlink{source:#1:#2}{[S#2]}}
\newcommand{\eref}[2]{\hyperlink{extra:#1:#2}{[E#2]}}
% Turn the next switch off for a shorter reading copy. The quotations preserve
% every exactTarget field, including qualifications omitted from a short summary.
\newif\ifshowcatalogue
\showcataloguetrue
\newcommand{\cataloguescope}[1]{\ifshowcatalogue
 \paragraph{Exact scope in the supplied catalogue.}
 \begin{quote}\small #1\end{quote}\fi}
% Compile twice with: lualatex 3349.tex
\numberwithin{equation}{section}
\hypersetup{pdftitle={Research attempt -- problem 3349},pdfauthor={Research notebook; original compilation retained}}
\fancyhead[L]{\small\sffamily Research notebook}
\fancyhead[R]{\small\sffamily Problem 3349 | 27 September 2026}
\begin{document}
\setcounter{section}{0}
% ================================================================
% problemId: 3349
% Canonical problem ID: 3349
% exactTarget SHA-256: a1ed443e0346c03af5174b78df861b2a10822c9aae61d5afd849732b1866b34b
\clearpage
\section*{\texorpdfstring{Algebraic independence of pi and e}{Algebraic independence of pi and e}}\label{3349}
\subsection{Definitions and mathematical statement}
Define $e=\sum_{n=0}^{\infty}1/n!$ and let $\pi$ be the positive half-period of the usual sine function. Algebraic independence of $e$ and $\pi$ over $\Q$ means
\[
 \forall P\in\Q[X,Y]\setminus\{0\},\qquad P(e,\pi)\ne0.
\]
Equivalently, the evaluation homomorphism $\Q[X,Y]\to\R$ is injective, or the transcendence degree of $\Q(e,\pi)$ over $\Q$ is two. This is stronger than showing either number individually transcendental or showing just $e+\pi$ irrational. A counterexample would require a nonzero polynomial relation with rational coefficients, not a numerical near-equality.
\par\smallskip\textit{Formulation and concept references:} \sref{3349}{1}.
\cataloguescope{Prove that no nonzero polynomial P in Q[X,Y] satisfies P(e,pi)=0.}
\subsection{Short English statement}
Are the constants e and pi free of every nontrivial polynomial relation with rational coefficients?
% BEGIN RESEARCH ADDITION ATTEMPT-168
\subsection{Research attempt}

\paragraph{Current frontier.}
Individual transcendence, algebraic independence of some other
exponential constants, and functional transcendence must be separated.
Pila's notes state the $e,\pi$ problem and explain Schanuel's conjecture
and functional analogues; the cited Nesterenko result involves
$\pi$ and $e^\pi$, not $e$ and $\pi$ \sref{3349}{1}, Sections 2--4.
A 2026 manuscript studies a symmetry between recurrences associated
with $e$ and $\pi$ and their meromorphic connection functions
\eref{3349}{1}, Sections 1 and 9. Such a relationship between functions
or limiting recurrences is not a polynomial relation between the two
numbers. No unconditional solution of the printed target was verified
in the current source search. The work below combines an algebraic
reformulation, an auxiliary-function attempt, and an exact finite
nonvanishing certificate.

\paragraph{Attack route.}
A conditional transcendence-degree argument is very short but imports
Schanuel's conjecture. A functional approach would need an arithmetic
specialization theorem that survives evaluation at the particular
arguments $1$ and $i\pi$. An auxiliary-function approach would need
integer linear forms in mixed monomials with simultaneously controlled
size, denominators, and nonvanishing. A search for a counterexample must
use exact certification rather than a numerical near-relation. Each
route is pursued far enough to expose its separate missing condition.

\paragraph{Attempt: symmetric coordinates lose no transcendence degree.}
Put $s=e+\pi$ and $t=e\pi$. Since $e$ and $\pi$ are roots of
$Z^2-sZ+t$, the field $\Q(e,\pi)$ is algebraic over $\Q(s,t)$, and
\begin{equation}\label{3349:symmetric}
 \operatorname{trdeg}_{\Q}\Q(e,\pi)
 =\operatorname{trdeg}_{\Q}\Q(s,t).
\end{equation}
More concretely, if a nonzero $P(X,Y)\in\Q[X,Y]$ vanished at $(e,\pi)$,
then $P(X,Y)P(Y,X)$ would be a nonzero symmetric polynomial. The
fundamental theorem of symmetric polynomials writes it uniquely as
$Q(X+Y,XY)$ for a nonzero $Q\in\Q[S,T]$. Thus every proposed relation
produces one in the symmetric coordinates, though possibly of higher
degree and height.

This does not show that both $s$ and $t$ are transcendental. It does
show that they cannot both be algebraic, since then $e$ and $\pi$ would
both be algebraic, contradicting their known individual transcendence
\sref{3349}{1}, Section 2. That elementary consequence is strictly weaker
than \eqref{3349:symmetric} having value two, and even an irrationality
proof for $s$ alone would not establish the target.

\paragraph{Attempt: the precise Schanuel dependency.}
The numbers $1$ and $i\pi$ are linearly independent over $\Q$, since
the real and imaginary parts of $a+bi\pi=0$ give $a=b=0$. Schanuel's
conjecture for just this pair would imply
\begin{equation}
 \operatorname{trdeg}_{\Q}\Q(1,i\pi,e,e^{i\pi})\geq2.
\end{equation}
Euler's identity gives $e^{i\pi}=-1$. Adjoining $i$ is algebraic and
does not change transcendence degree, so the conclusion would be
$\operatorname{trdeg}_{\Q}\Q(e,\pi)=2$. This is the familiar conditional
route in the framework of \sref{3349}{1}, Section 3, not a proof of the
needed numerical instance of Schanuel. The full target also implies
transcendence of both $e+\pi$ and $e\pi$, by substituting $X+Y$ or $XY$
into a proposed algebraic equation. These consequences are useful
warnings against treating a new ``specialization lemma'' as innocuous.

\paragraph{Attempt: why functional independence does not specialize.}
For an irrational real $\alpha$, $\exp z$ and $\exp(\alpha z)$ are
algebraically independent as analytic functions over $\C$. A polynomial
relation would be a finite linear combination of
$\exp((m+n\alpha)z)$ with distinct exponents. Differentiating at zero
produces a Vandermonde system, so every coefficient vanishes. Yet at
$z=0$ both functions equal one. This proves, without any transcendence
conjecture, that even exact functional algebraic independence can have
an algebraically dependent special value.

The functional results discussed in \sref{3349}{1}, Sections 4--6,
control positive-dimensional analytic or algebraic configurations.
Applying one directly to the isolated pair of arguments $(1,i\pi)$
would require an additional arithmetic statement; the functional
result alone does not supply it. The same problem affects a proposed
use of the recurrence symmetries in \eref{3349}{1}: relations among
connection functions have to be accompanied by a theorem about these
particular values, not just their functional equations.

\paragraph{Attempt: a mixed auxiliary-function calculation.}
Fix an integer $d\geq1$ and seek
\begin{equation}
 F(z)=\sum_{0\leq i,j\leq d}c_{ij}z^i e^{jz},
 \qquad c_{ij}\in\Z,
\end{equation}
with a large zero at $z=0$. There are $M=(d+1)^2$ coefficients.
Vanishing of the first $M-1$ Taylor coefficients gives $M-1$ rational
homogeneous linear equations, so a nonzero rational solution exists;
clearing denominators gives an integral one. The functions $z^ie^{jz}$
are linearly independent: on the positive real axis, divide a putative
identity by the largest exponential and let $z\to+\infty$; its leading
polynomial would tend to zero, so must vanish. Induction removes all
exponentials. Thus this construction gives a genuinely nonzero $F$.

Its two evaluations have sharply different arithmetic forms:
\begin{equation}\label{3349:endpoints}
 F(1)=\sum_{j=0}^d\Bigl(\sum_i c_{ij}\Bigr)e^j,
 \qquad
 F(i\pi)=\sum_{i=0}^d\Bigl(\sum_j(-1)^jc_{ij}\Bigr)(i\pi)^i.
\end{equation}
The second expression uses exactly $e^{ji\pi}=(-1)^j$. In this
maximal-order construction one can in fact prove that \emph{both}
collected polynomials in \eqref{3349:endpoints} are nonzero. Here is a
useful check that goes beyond merely observing $F\not\equiv0$.
An exponential polynomial of degrees $i\le u$, $j\le v$ is annihilated
by the constant-coefficient differential operator
\begin{equation}
 \prod_{j=0}^v\left(\frac{\mathrm d}{\mathrm dz}-j\right)^{u+1},
\end{equation}
of order $(u+1)(v+1)$. If its first that many derivatives vanish at
zero, the differential equation recursively forces every derivative
to vanish, so the function is identically zero.

Write $P_0(Z,W)=\sum c_{ij}Z^iW^j$. If the first polynomial in
\eqref{3349:endpoints} were identically zero, then $P_0(1,W)=0$ and
$P_0$ would be divisible by $Z-1$. Dividing $F(z)$ by $z-1$ would give
an exponential polynomial with degree bounds $(d-1,d)$ and the same
zero of order at least $M-1$ at zero. Its annihilating operator has
order $d(d+1)\le M-1$, an impossibility. If the second polynomial were
identically zero, then $P_0(Z,-1)=0$ and $P_0$ would be divisible by
$W+1$. Dividing $F(z)$ by $e^z+1$, which is nonzero at zero, similarly
gives degree bounds $(d,d-1)$ and the same contradiction. Thus both
endpoint polynomials are nonzero; individual transcendence of $e$ and
$\pi$ now proves $F(1)F(i\pi)\ne0$ \sref{3349}{1}, Section 2.
This is nonvanishing of these \emph{separated} forms, not algebraic
independence of all mixed monomials. The script independently checks
the exact rational interpolation matrix, zero order, and nonzero
collected endpoint polynomials for $1\le d\le4$; the argument above
covers every $d$.

There is also an explicit coefficient bound. Use derivative equations
rather than Taylor coefficients: their matrix entries are
$(r!/(r-i)!)j^{r-i}$ for $r\ge i$, and zero otherwise, for
$0\le r\le M-2$, with $0^0=1$. Its rank is $M-1$ by the same
annihilating-operator argument. Taking maximal minors gives a nonzero
integer kernel vector satisfying, for $d\ge2$,
\begin{equation}
 H:=\max|c_{ij}|
 \le\left(\sqrt{M-1}\,M^d d^{M-2}\right)^{M-1}.
\end{equation}
This is Hadamard's inequality applied to the minor matrices. If
$m=M-1$ and $\rho>|z|$, the maximum principle for $F(z)/z^m$ gives
\begin{equation}
 |F(z)|\le M H\max(1,\rho^d)e^{d\rho}
                   \left(\frac{|z|}{\rho}\right)^m.
\end{equation}
Choosing $\rho=m/d$ makes the negative logarithmic contribution at a
fixed endpoint of scale $-d^2\log d$, but the available upper bound on
$\log H$ has scale $d^4\log d$. Substituting this coefficient bound
therefore does not even force these endpoint values to be small. This
failure is about the bounds proved here; it does not show that sharper
auxiliary functions or sharper coefficient estimates are impossible.

To derive a contradiction from an arbitrary $P(e,\pi)=0$, one still
needs an arithmetic lower bound for suitable mixed forms, with
coefficient and denominator growth dominated by analytic smallness.
The separated nonvanishing lemma supplies no such mixed estimate.
If maximal-order interpolation is discarded, the naive nonvanishing
principle fails even more directly: $(e^z-1)^m(e^z+1)$ is nonzero as a
function, vanishes to order $m$ at zero, and vanishes at $i\pi$.
Thus neither a Taylor remainder nor high-order vanishing by itself is
an arithmetic independence theorem.

\paragraph{Exact computation: a bounded polynomial family.}
The accompanying standard-library Python script uses rational interval
arithmetic, not floating-point root recognition. It proves the finite
statement
\begin{equation}\label{3349:finite}
 \left|P(e,\pi)\right|>2^{-10}
 \quad\begin{gathered}
  \text{for every nonzero }P\in\Z[X,Y]\text{ of total degree at most }3,\\
  \text{with every coefficient in }\{-1,0,1\}.
 \end{gathered}
\end{equation}
There are $3^{10}-1=59\,048$ such polynomials. Set
$E_N=\sum_{j=0}^N1/j!$. A geometric-series comparison gives
\begin{equation}
 E_N<e<E_N+\frac1{N\,N!}\qquad(N\geq1).
\end{equation}
For $0<x<1$, consecutive even and odd truncations of
$\arctan x=\sum_{j\geq0}(-1)^j x^{2j+1}/(2j+1)$ enclose its value.
Machin's identity
\begin{equation}
 \pi=16\arctan(1/5)-4\arctan(1/239)
\end{equation}
follows by the tangent addition formula and the fact that
$4\arctan(1/5)-\arctan(1/239)$ lies in $(0,\pi/2)$ with tangent one.
The factorial sum through $1/80!$ and eighty arctangent terms suffice for the
following dyadic outer enclosures. If $D=2^{128}$, then
\begin{align*}
 L_e&=924983374546220337150911035843336795079,
 & L_e&<De<L_e+1,\\
 L_\pi&=1069028584064966747859680373161870783300,
 & L_\pi&<D\pi<L_\pi+1.
\end{align*}
For readability the script, rather than a table of decimal values,
records the exact integers and the assertion that the analytic
intervals lie inside these outer bounds.

All monomials are positive at the endpoints. Bound $e^i\pi^j$ by the
products of the lower and upper endpoint powers, and reverse endpoints
when a coefficient is negative. Put every result over the common
denominator $D^3$, including terms of smaller degree. Exhaustive
integer addition verifies that each polynomial interval lies entirely
outside $[-2^{-10},2^{-10}]$. The polynomial attaining the smallest
certified distance in this enumeration is
\begin{equation}
 P_*(X,Y)=X^3-Y^3+X^2-Y^2+XY+X+Y-1
\end{equation}
or its negative. Its closeness to zero is not a relation; the interval
proof explicitly separates it from zero. This is a finite, reproducible
nonvanishing proof for \eqref{3349:finite}, not a probabilistic search
claim. No assertion is made for other degrees or coefficient heights.

\paragraph{Stress test and outcome.}
\textbf{Outcome: Exploratory approach with unresolved gap.}

The symmetric reformulation is exact, the Schanuel implication is
explicitly conditional, and the auxiliary-function calculation proves separated endpoint
nonvanishing but stops at an uncontrolled mixed arithmetic estimate. The computation proves
a bounded-family certificate with rational error control; it does not
promote a finite search or a very small decimal to an algebraic
independence theorem. The universal quantifier over all nonzero rational
polynomials, equivalently all integral polynomials after clearing
denominators, remains unhandled. No counterexample polynomial or
unconditional complete proof candidate has been obtained.
% END RESEARCH ADDITION ATTEMPT-168
\subsection{Sources}
\begin{itemize}\raggedright
\item[\textnormal{[S1]}]\hypertarget{source:3349:1}{} Jonathan Pila, Functional transcendence via o-minimality, edited LMS-EPSRC 2013 lecture notes.
\url{https://people.maths.ox.ac.uk/pila/LMSNotes.pdf}.
{\small\textit{Catalogue formulation source.}}
% BEGIN RESEARCH ADDITION SOURCES-168
\item[\textnormal{[E1]}]\hypertarget{extra:3349:1}{} Benoit Cloitre, ``The holonomic triangle: from a symmetry between $e$ and $\pi$ to additive Gamma functions,'' arXiv:2601.04242v1 (January 2026). Sections 1 and 9.
\url{https://arxiv.org/pdf/2601.04242}.
{\small\textit{Recent recurrence/functional connection, not a polynomial relation or an independence proof. Only the stated scope is used.}}
% END RESEARCH ADDITION SOURCES-168
\end{itemize}

\end{document}
